\setcounter{section}{1}
\setcounter{theorem}{0}
\setcounter{definition}{0}
\setcounter{remark}{0}
\setcounter{note}{0}

\subsubsection{Preliminaries: Euler--Maclaurin, after Qingshu}

The official CAM syllabus names this circle as interpolation and approximation, including trigonometric interpolation \cite{conte00}. Problem~1 is the periodic trapezoidal rule. Part~(i) is the statement that periodicity upgrades the usual \(O(h^{2})\) error to the full smoothness order \(O(h^{m})\). Part~(ii) is what that becomes when \(f\) is analytic: the error is exponential in \(n=2\pi/h\).

The piecewise-linear interpolation remainder does \emph{not} give (i). On each panel \([ih,(i+1)h]\) one has
\[
f(x)-p_{1}(x)
=
\frac{f''(\xi_{x})}{2}(x-ih)\bigl(x-(i+1)h\bigr),
\]
and integrating produces the classical composite estimate \(I(f)-I_{h}^{Tr}(f)=-\frac{2\pi}{12}h^{2}f''(\xi)=O(h^{2})\). Higher smoothness only improves the constant. The missing mechanism is cancellation of the endpoint corrections, which is exactly Euler--Maclaurin. The background below is Qingshu, \S8.5 (printed pp.~94--97; PDF pp.~102--107) \cite{qingshu11}.

\paragraph{Qiuzhen's quadrature questions, 2025--2026.}

{\footnotesize
\begin{center}
\begin{tabular}{@{}>{\raggedright\arraybackslash}p{2.45cm}
                   >{\raggedright\arraybackslash}p{4.35cm}
                   >{\raggedright\arraybackslash}p{3.45cm}
                   >{\raggedright\arraybackslash}p{3.55cm}@{}}
\toprule
Paper & Typical ask & Machine & Relation to 2025 Spring \\
\midrule
2025 Spring P1 & periodic trapezoid, order \(m\) / exponential & Euler--Maclaurin; Fourier aliasing & this problem \\
2025 Fall P2 & Gauss--Chebyshev nodes and weights & interpolatory quadrature at zeros of \(T_{n+1}\) & same word ``quadrature'', different nodes \\
2026 Spring P2 & Gauss--Legendre error via \(\omega^{2}\) & Hermite remainder / Peano & Qingshu Ch.~10, not Ch.~8 \\
\bottomrule
\end{tabular}
\end{center}
}

\begin{definition}[Periodic \(C^{m}\), and the two trapezoidal sums]
\label{def:p1-periodic}
A \(2\pi\)-periodic function of class \(C^{m}\) is a function \(f\in C^{m}(\mathbb{R})\) with \(f(x+2\pi)=f(x)\). Equivalently, \(f\in C^{m}([0,2\pi])\) and
\[
f^{(j)}(0)
=
f^{(j)}(2\pi),
\qquad
j=0,1,\dots,m.
\]
(The exam's phrase ``periodic'' \(+\) \(f\in C^{m}([0,2\pi])\) is read this way: without matching jets at the endpoints the periodic extension is not \(C^{m}\), and the order in (i) drops to the first unmatched derivative.) Write
\[
I(f)
\coloneqq
\int_{0}^{2\pi} f(x)\,\mathrm{d}x,
\qquad
h
=
\frac{2\pi}{n},
\qquad
I_{h}^{Tr}(f)
\coloneqq
h\sum_{i=1}^{n} f(ih).
\]
The ordinary composite trapezoidal sum on \([0,2\pi]\) is
\begin{equation}
\label{eq:p1-ordinary-trap}
T_{h}(f)
\coloneqq
\frac{h}{2}\bigl(f(0)+f(2\pi)\bigr)
+
h\sum_{i=1}^{n-1} f(ih).
\end{equation}
If \(f(0)=f(2\pi)\), then \(I_{h}^{Tr}(f)=T_{h}(f)\). The exam writes the periodic form; the Euler--Maclaurin identity is most cleanly stated for \(T_{h}\).
\end{definition}

\paragraph{Qingshu's periodic Bernoulli functions.}

Qingshu does not start from the classical Bernoulli numbers. The functions \(b_{n}\) are defined by integrating the sawtooth and imposing mean zero; the values \(b_{n}(0)\) then become the coefficients in Euler--Maclaurin.

\begin{definition}[Periodic functions \(b_{n}\); Qingshu, Definition~8.1]
\label{def:p1-bn}
Let \(b_{1}\) be the \(1\)-periodic extension of \(x-\tfrac12\) from \([0,1)\). For each \(n\ge 2\), let \(b_{n}\) be the unique \(1\)-periodic antiderivative of \(b_{n-1}\) with
\[
\int_{0}^{1} b_{n}(x)\,\mathrm{d}x
=
0.
\]
Thus \(b_{n}'=b_{n-1}\) almost everywhere, and \(b_{n}(x)=B_{n}(\{x\})/n!\) in the usual Bernoulli-polynomial notation. On \([0,1)\) the first four are
\begin{align*}
b_{1}(x)
&=
x-\tfrac12,
&
b_{2}(x)
&=
\tfrac12 x^{2}-\tfrac12 x+\tfrac{1}{12},
\\
b_{3}(x)
&=
\tfrac16 x^{3}-\tfrac14 x^{2}+\tfrac{1}{12}x,
&
b_{4}(x)
&=
\tfrac{1}{24}x^{4}-\tfrac{1}{12}x^{3}+\tfrac{1}{24}x^{2}-\tfrac{1}{720}.
\end{align*}
In particular \(b_{2}(0)=\frac{1}{12}\) and \(b_{4}(0)=-\frac{1}{720}\).
\end{definition}

\begin{proposition}[Fourier series of \(b_{n}\); Qingshu, Proposition~8.4]
\label{prop:p1-bn-fourier}
For \(k\in\mathbb{N}\),
\begin{align*}
b_{1}(x)
&\sim
-\frac{1}{\pi}\sum_{n=1}^{\infty}\frac{\sin(2n\pi x)}{n},
\\
b_{2k}(x)
&=
\frac{2(-1)^{k+1}}{(2\pi)^{2k}}\sum_{n=1}^{\infty}\frac{\cos(2n\pi x)}{n^{2k}},
\\
b_{2k+1}(x)
&=
\frac{2(-1)^{k+1}}{(2\pi)^{2k+1}}\sum_{n=1}^{\infty}\frac{\sin(2n\pi x)}{n^{2k+1}}.
\end{align*}
The series for \(b_{1}\) converges to \(b_{1}(x)\) off \(\mathbb{Z}\) and to \(0\) at integers (Dirichlet).
\end{proposition}

\begin{corollary}[Values and bounds; Qingshu, Corollary~8.2]
\label{cor:p1-bn-bound}
For every \(k\in\mathbb{N}\),
\[
b_{2k+1}(0)
=
0,
\qquad
\lvert b_{k}(x)\rvert
\le
\frac{\zeta(k)}{(2\pi)^{k}}\qquad (k\ge 2),
\]
and \(\lvert b_{1}(x)\rvert\le\tfrac12\).
\end{corollary}

\paragraph{Qingshu's Euler--Maclaurin formula.}

The unit-step identity is the one printed in the notes. The exam needs the same identity after the change of scale \(x=th\).

\begin{proposition}[Order-zero formula; Qingshu, Proposition~8.3 / (8.41)]
\label{prop:p1-em0}
Let \(a,b\in\mathbb{Z}\) and \(f\in C^{1}[a,b]\). Then
\begin{equation}
\label{eq:p1-em0}
\sum_{k=a}^{b} f(k)
=
\int_{a}^{b} f(x)\,\mathrm{d}x
+
\frac{f(a)+f(b)}{2}
+
\int_{a}^{b} b_{1}(x)\,f'(x)\,\mathrm{d}x.
\end{equation}
\end{proposition}

\begin{theorem}[Euler--Maclaurin; Qingshu, Theorem~8.8 / (8.42)]
\label{thm:p1-em}
Let \(a,b,m\in\mathbb{Z}\) with \(m\ge 2\), and let \(f\in C^{m}[a,b]\). Then
\begin{equation}
\label{eq:p1-em}
\sum_{k=a}^{b} f(k)
-
\int_{a}^{b} f(x)\,\mathrm{d}x
=
\frac{f(a)+f(b)}{2}
+
\sum_{k=2}^{m}
\bigl[f^{(k-1)}(b)-f^{(k-1)}(a)\bigr]\,b_{k}(0)
+
(-1)^{m+1}
\int_{a}^{b} b_{m}(x)\,f^{(m)}(x)\,\mathrm{d}x.
\end{equation}
\end{theorem}

\begin{proof}[Proof, after Qingshu]
The case \(m=1\) is \eqref{eq:p1-em0}. For \(m=2\), integrate by parts using \(b_{2}'=b_{1}\) and the \(1\)-periodicity of \(b_{2}\) (so \(b_{2}(b)=b_{2}(a)=b_{2}(0)\)):
\[
\int_{a}^{b} b_{1}(x)\,f'(x)\,\mathrm{d}x
=
\bigl[f'(b)-f'(a)\bigr]\,b_{2}(0)
-
\int_{a}^{b} b_{2}(x)\,f''(x)\,\mathrm{d}x.
\]
Substitute into \eqref{eq:p1-em0}. The general formula is the same integration by parts, repeated.
\end{proof}

\begin{corollary}[Even form; Qingshu, Corollary~8.3 / (8.44)]
\label{cor:p1-em-even}
If \(f\in C^{2m}[a,b]\), then odd endpoint terms drop by \(b_{2k+1}(0)=0\), and
\begin{equation}
\label{eq:p1-em-even}
\sum_{k=a}^{b} f(k)
-
\int_{a}^{b} f(x)\,\mathrm{d}x
=
\frac{f(a)+f(b)}{2}
+
\sum_{k=1}^{m}
\bigl[f^{(2k-1)}(b)-f^{(2k-1)}(a)\bigr]\,b_{2k}(0)
-
\int_{a}^{b} b_{2m}(x)\,f^{(2m)}(x)\,\mathrm{d}x.
\end{equation}
\end{corollary}

\paragraph{How to use it on the exam interval.}

Qingshu sums \(f(k)\) with step \(1\). The exam samples \(f(ih)\) with step \(h\). The reduction is the substitution \(F(t)=f(th)\).

\begin{lemma}[Euler--Maclaurin at step \(h\)]
\label{lem:p1-em-h}
Let \(h=2\pi/n\) and \(f\in C^{m}([0,2\pi])\), with \(m\ge 1\). Then
\begin{equation}
\label{eq:p1-em-h}
T_{h}(f)-I(f)
=
\sum_{k=2}^{m}
h^{k}
\bigl[f^{(k-1)}(2\pi)-f^{(k-1)}(0)\bigr]\,b_{k}(0)
+
R_{m}(f,h),
\end{equation}
where the \(m=1\) sum is empty, and
\begin{equation}
\label{eq:p1-Rm}
R_{m}(f,h)
=
(-1)^{m+1}
h^{m}
\int_{0}^{2\pi}
b_{m}\!\left(\frac{x}{h}\right)
f^{(m)}(x)\,
\mathrm{d}x.
\end{equation}
In particular
\begin{equation}
\label{eq:p1-Rm-bound}
\bigl\lvert R_{m}(f,h)\bigr\rvert
\le
C_{m}\,h^{m}\,\lVert f^{(m)}\rVert_{L^{1}[0,2\pi]},
\qquad
C_{1}
=
\tfrac12,
\qquad
C_{m}
=
\frac{\zeta(m)}{(2\pi)^{m}}\quad (m\ge 2).
\end{equation}
\end{lemma}

\begin{proof}
Set \(F(t)\coloneqq f(th)\) for \(t\in[0,n]\), so \(F^{(j)}(t)=h^{j}f^{(j)}(th)\). Apply \cref{prop:p1-em0} if \(m=1\), and \cref{thm:p1-em} if \(m\ge 2\), on the integers from \(0\) to \(n\):
\begin{align*}
\sum_{k=0}^{n} F(k)
-
\int_{0}^{n} F(t)\,\mathrm{d}t
&=
\frac{F(0)+F(n)}{2}
+
\sum_{k=2}^{m}
\bigl[F^{(k-1)}(n)-F^{(k-1)}(0)\bigr]\,b_{k}(0)
\\
&\qquad
+
(-1)^{m+1}
\int_{0}^{n} b_{m}(t)\,F^{(m)}(t)\,\mathrm{d}t.
\end{align*}
The integral on the left is \(h^{-1}I(f)\). Multiply through by \(h\). The left-hand side becomes
\[
h\sum_{k=0}^{n} f(kh)
-
I(f)
=
T_{h}(f)
+
\frac{h}{2}\bigl(f(0)+f(2\pi)\bigr)
-
I(f),
\]
and the displayed endpoint sum on the right is exactly \(\frac{h}{2}(f(0)+f(2\pi))\) plus the Bernoulli corrections. What remains is \(T_{h}(f)-I(f)\) as in \eqref{eq:p1-em-h}. The remainder integral is
\[
(-1)^{m+1}h^{m+1}
\int_{0}^{n}
b_{m}(t)\,f^{(m)}(th)\,\mathrm{d}t.
\]
The substitution \(x=th\) produces \eqref{eq:p1-Rm}. The bound \eqref{eq:p1-Rm-bound} is \cref{cor:p1-bn-bound}.
\end{proof}

\begin{remark}[What periodicity does]
\label{rem:p1-cancel}
Every coefficient in the sum \eqref{eq:p1-em-h} is a jump \(f^{(k-1)}(2\pi)-f^{(k-1)}(0)\). On a non-periodic interval those jumps are the reason the trapezoidal rule stays \(O(h^{2})\) no matter how smooth \(f\) is (the first surviving term is \(h^{2}b_{2}(0)(f'(2\pi)-f'(0))\)). If \(f\) is \(2\pi\)-periodic of class \(C^{m}\), every jump through order \(m-1\) vanishes, and \eqref{eq:p1-em-h} collapses to \(T_{h}(f)-I(f)=R_{m}(f,h)=O(h^{m})\). That is the whole of part~(i).
\end{remark}

\subsubsection{Solution of (i)}

\begin{theorem}[\(C^{m}\) periodic \(\Rightarrow\) order \(m\)]
\label{thm:p1-order-m}
Let \(f\) be \(2\pi\)-periodic of class \(C^{m}\), \(m\ge 1\). Then \(I_{h}^{Tr}(f)=T_{h}(f)\) and
\[
\bigl\lvert I(f)-I_{h}^{Tr}(f)\bigr\rvert
\le
C_{m}\,h^{m}\,\lVert f^{(m)}\rVert_{L^{1}[0,2\pi]},
\]
with \(C_{m}\) as in \eqref{eq:p1-Rm-bound}. In particular the exam bound holds with \(C=C_{m}\lVert f^{(m)}\rVert_{L^{1}}\).
\end{theorem}

\begin{proof}
Periodicity gives \(f(0)=f(2\pi)\), hence \(I_{h}^{Tr}=T_{h}\) by \cref{def:p1-periodic}. The same periodicity gives \(f^{(j)}(0)=f^{(j)}(2\pi)\) for \(0\le j\le m\), so every Bernoulli endpoint term in \cref{lem:p1-em-h} vanishes. The identity reduces to \(I_{h}^{Tr}(f)-I(f)=R_{m}(f,h)\), and \eqref{eq:p1-Rm-bound} is the claim.
\end{proof}

\begin{remark}[Fourier aliasing, and why it is not the proof of (i) for \(m=1\)]
\label{rem:p1-alias-cm}
Write \(\hat f(k)=(2\pi)^{-1}\int_{0}^{2\pi}f(x)e^{-ikx}\,\mathrm{d}x\). If the Fourier series of \(f\) may be summed absolutely,
\begin{equation}
\label{eq:p1-alias}
I(f)
=
2\pi\,\hat f(0),
\qquad
I_{h}^{Tr}(f)
=
2\pi\sum_{\ell\in\mathbb{Z}}\hat f(\ell n),
\end{equation}
hence
\[
I(f)-I_{h}^{Tr}(f)
=
-2\pi\sum_{\ell\neq 0}\hat f(\ell n).
\]
(The inner geometric sum \(h\sum_{j=1}^{n}e^{ikjh}\) equals \(2\pi\) if \(n\mid k\) and equals \(0\) otherwise.) Integration by parts on the circle gives \(\lvert\hat f(k)\rvert\le (2\pi)^{-1}\lvert k\rvert^{-m}\lVert f^{(m)}\rVert_{L^{1}}\). For \(m\ge 2\) one may sum the aliases:
\[
\bigl\lvert I(f)-I_{h}^{Tr}(f)\bigr\rvert
\le
n^{-m}\,\lVert f^{(m)}\rVert_{L^{1}}\cdot 2\zeta(m)
=
O(h^{m}).
\]
For \(m=1\) the comparison series \(\sum_{\ell\neq 0}\lvert\ell\rvert^{-1}\) diverges, so this bound does not close. Euler--Maclaurin is the argument that covers every \(m\ge 1\), including the exam's \(C^{1}\) case.
\end{remark}

\subsubsection{Solution of (ii)}

``Analytic'' for a \(2\pi\)-periodic function means the periodic extension is real-analytic on \(\mathbb{R}\). Equivalently, \(f\) extends holomorphically to a strip \(\lvert\operatorname{Im}z\rvert<\alpha\) for some \(\alpha>0\).

\begin{lemma}[Exponential decay of Fourier coefficients]
\label{lem:p1-fourier-exp}
If \(f\) is \(2\pi\)-periodic and holomorphic in \(\lvert\operatorname{Im}z\rvert<\alpha\), then for every \(\beta\in(0,\alpha)\) there is \(M_{\beta}<\infty\) such that
\[
\lvert\hat f(k)\rvert
\le
M_{\beta}\,e^{-\beta\lvert k\rvert},
\qquad
k\in\mathbb{Z}.
\]
\end{lemma}

\begin{proof}
Shift the integration contour of \(\hat f(k)\) to \(\operatorname{Im}z=\beta\operatorname{sign}(-k)\) if \(k\neq 0\) (justified by periodicity in the real direction and holomorphy in the strip). Then \(\lvert e^{-ikz}\rvert=e^{-\beta\lvert k\rvert}\), and \(M_{\beta}\) is \((2\pi)^{-1}\) times the maximum of \(\lvert f\rvert\) on the two lines \(\operatorname{Im}z=\pm\beta\), times the length \(2\pi\).
\end{proof}

\begin{theorem}[Analytic periodic \(\Rightarrow\) exponential convergence]
\label{thm:p1-exp}
Let \(f\) be \(2\pi\)-periodic and holomorphic in a strip of width \(\alpha>0\). Then for every \(\beta\in(0,\alpha)\) there is \(C_{\beta}<\infty\) such that
\[
\bigl\lvert I(f)-I_{h}^{Tr}(f)\bigr\rvert
\le
C_{\beta}\,e^{-\beta n}
=
C_{\beta}\,e^{-2\pi\beta/h}.
\]
\end{theorem}

\begin{proof}
Absolute convergence of the Fourier series is \cref{lem:p1-fourier-exp}, so the aliasing identity \eqref{eq:p1-alias} applies and
\[
\bigl\lvert I(f)-I_{h}^{Tr}(f)\bigr\rvert
\le
2\pi\sum_{\ell\neq 0}
M_{\beta}\,e^{-\beta n\lvert\ell\rvert}
=
\frac{4\pi M_{\beta}\,e^{-\beta n}}{1-e^{-\beta n}}.
\]
The right-hand side is \(O(e^{-\beta n})\).
\end{proof}

\begin{remark}[The same conclusion from Euler--Maclaurin alone]
\label{rem:p1-em-exp}
If one prefers not to name Fourier coefficients, apply \cref{thm:p1-order-m} at every order. Cauchy's estimate on a circle of radius \(r\in(0,\alpha)\) in the complex plane gives \(\lVert f^{(m)}\rVert_{\infty}\le M\,m!\,r^{-m}\). Combined with \eqref{eq:p1-Rm-bound},
\[
\bigl\lvert I(f)-I_{h}^{Tr}(f)\bigr\rvert
\le
C\,\Bigl(\frac{h}{2\pi r}\Bigr)^{m}m!.
\]
Stirling's bound \(m!\le (m/e)^{m}\sqrt{2\pi m}\,e^{1/(12m)}\) and the choice \(m\sim 2\pi r/h\) produce an estimate \(O(e^{-c/h})\) for some \(c>0\). This is the same exponential, read off Qingshu's remainder rather than off aliasing. The Fourier argument is shorter.
\end{remark}

\begin{examnote}[What the two parts are]
\label{examnote:p1}
\begin{enumerate}[label=\textup{(\roman*)}]
    \item Quote Qingshu's formula at step \(h\), \cref{lem:p1-em-h}. Periodicity kills every jump \(f^{(k-1)}(2\pi)-f^{(k-1)}(0)\). The error is the integral remainder \(R_{m}=O(h^{m})\). Do not stop at the piecewise-linear interpolation remainder: that is only \(O(h^{2})\).
    \item Analytic \(\Rightarrow\) holomorphic in a strip \(\Rightarrow\) \(\lvert\hat f(k)\rvert\le M e^{-\beta\lvert k\rvert}\). The trapezoidal sum aliases to \(2\pi\sum_{\ell}\hat f(\ell n)\), so the error is \(O(e^{-\beta n})\). The Euler--Maclaurin-only route is \cref{rem:p1-em-exp}.
\end{enumerate}
\end{examnote}

\begin{remark}[Local Qingshu slice]
\label{rem:p1-qingshu-em}
The statements copied above are Qingshu, \S8.5.1, Theorem~8.8 and Corollary~8.3 \cite{qingshu11}. A six-page excerpt (original PDF pp.~102--107) sits at \texttt{qe-review/applied-math/reference/qingshu-euler-maclaurin/}.
\end{remark}

% Restore parent section so subsequent \subsection headings stay under
% ``Statement and solutions''.
\setcounter{section}{2}
